\begin{lstlisting}[
    caption={Viewpoint Adaptation}, 
    label={alg:graph-combination},
    basicstyle=\ttfamily\small,
    breaklines=true,
    frame=single,
    numbers=none,
    keepspaces=true,
    showstringspaces=false,
    mathescape=true
]
Algorithm: InstantiateSummary(G_caller, summary_callee, actualArgs, returnVar)

Input:  Caller's current graph G_caller,
        callee's exit-graph summary G_callee = (I_c, O_c, L_c, E_c, W_c),
        actual arguments [a0, ..., ak],
        variable to receive return value (may be null)
Output: G_caller updated in place

// --- Step 0: Remap callee node IDs to fresh caller-namespace IDs ---
//     (InsideNodes and LoadNodes get fresh IDs; ParameterNodes are not remapped;
//      GlobalNode maps to itself)

for each InsideNode n in G_callee do
    remap(n) $\leftarrow$ fresh InsideNode in G_caller
for each LoadNode n in G_callee do
    remap(n) $\leftarrow$ fresh LoadNode in G_caller
Apply remap to all edges, return targets, E_c, and W_c in G_callee

// --- Step 1: Compute node mapping mu (least fixed point) ---

Initialize mu:
    for i = 0, ..., k do
        mu(P_i) $\leftarrow$ pt_caller(actualArgs[i])      // Constraint 1
    for all other callee nodes n do
        mu(n) $\leftarrow$ {}

repeat until mu reach a fixed-point:
    // Constraint 2: callee outside edge meets caller inside edge
    for each outside edge (n1 --f--> n2) in O_c do
        for each node n3 in mu(n1) do
            for each node n4 in InsideEdgeTargets_caller(n3, f) do
                mu(n2) $\leftarrow$ mu(n2) $\cup$ { n4 }

    // Constraint 3: callee outside edge meets callee inside edge (aliasing)
    for each outside edge (n1 --f--> n2) in O_c do
        for each inside edge (n3 --f--> n4) in I_c do
            S1 $\leftarrow$ mu(n1) $\cup$ ({n1} \ ParameterNode)
            S3 $\leftarrow$ mu(n3) $\cup$ ({n3} \ ParameterNode)
            if S1 intersects S3 and (n1 != n3 or n1 is LoadNode) then
                mu(n2) $\leftarrow$ mu(n2) $\cup$ mu(n4) $\cup$ ({n4} \ ParameterNode)

// Extended mapping: identity for non-parameter nodes
mu'(n) = mu(n) $\cup$ ({n} \ ParameterNode)

// --- Step 2 and 3: Combine graphs ---

// Inside edges
for each inside edge (n1 --f--> n2) in I_c do
    for each n1' in mu'(n1), n2' in mu'(n2) do
        addInsideEdge_caller(n1', f, n2')

// Outside edges (skip if source maps to InsideNode)
for each outside edge (n --f--> nL) in O_c do
    for each n' in mu'(n) do
        if n' is not InsideNode then
            addOutsideEdge_caller(n', f, nL)

// Return value
if returnVar != null then
    retNodes $\leftarrow$ Union of mu'(n) for each n in returnTargets(G_callee)
    strongUpdate(returnVar, retNodes)

// Escaped nodes
E_caller $\leftarrow$ E_caller $\cup$ { n' : n in E_c, n' in mu'(n) }

// Update Write Set
for each (n, f) in W_c do
    for each n' in mu'(n) do
        if n' is not InsideNode and n' exists in G_caller then
            recordMutation_caller(n', f)

// --- Step 4: Remove captured load nodes ---
//     A LoadNode is "captured" if it is unreachable from any root
//     (parameters, global escaped nodes, variable targets).
liveNodes $\leftarrow$ BFS from roots of G_caller along all edges
for each LoadNode nL in G_caller do
    if nL not in liveNodes then
        removeNode(nL)
\end{lstlisting}